\chapter{Введение}
\label{ch:intro}

В энергетике и холодильной технике точный расчёт тепловых процессов критически важен для проектирования эффективных систем. Особую сложность представляет описание поведения газов при сжатии и расширении, где ключевую роль играет разница между изобарной $C_p$ и изохорной $C_v$ теплоёмкостями. Неверные оценки этих величин могут привести к перерасходу энергии, снижению КПД оборудования или даже повреждению установок из-за перегрева.

Первый способ позволяет избежать сложных расчётов. Он предлагает считать теплоёмкость постоянной величиной в рабочем диапазоне температур. Однако такой подход даёт приемлемую точность только для узких интервалов условий. В реальных системах, где газы подвергаются значительному сжатию или резкому расширению (например, в турбинах или компрессорах), это приводит к существенным ошибкам в прогнозировании рабочих параметров.

Оптимальное решение подразумевает использование точного отношения $C_p/C_v$, которое позволяет корректно описывать термодинамические процессы в газах. Это особенно важно при расчетах адиабатических изменений состояния рабочего тела в тепловых машинах. Учёт этого отношения обеспечивает точное моделирование реальных процессов и помогает создавать энергоэффективные технические системы.

Целью работы являлось определение показателя адиабаты углекислого газа методом Клемана-Дезорма.